% from ega0-1-fr.tex
\oldpage[0\textsubscript{I}]{49}

et lorsque nous parlerons de faisceaux localement libres sans préciser, il
sera sous-entendu qu'ils sont de \emph{rang fini}.
\end{env}

\begin{env}[5.4.2]
\label{0.5.4.2-fr}
Si $\mathcal{L}$, $\mathcal{F}$ sont deux $\mathcal{O}_X$-Modules, on a un
homomorphisme canonique fonctoriel
\[
  \check{\mathcal{L}}\otimes_{\mathcal{O}_X}\mathcal{F}
  =\mathcal{H}\!om_{\mathcal{O}_X}(\mathcal{L},\mathcal{O}_X)
   \otimes_{\mathcal{O}_X}\mathcal{F}
  \longrightarrow
  \mathcal{H}\!om_{\mathcal{O}_X}(\mathcal{L},\mathcal{F})
  \tag{5.4.2.1}\label{0.5.4.2.1-fr}
\]
défini de la façon suivante~: pour tout ouvert $U$, à tout couple $(u,t)$, où
\[
  u\in\Gamma\bigl(U,
    \mathcal{H}\!om_{\mathcal{O}_X}(\mathcal{L},\mathcal{O}_X)\bigr)
  =\operatorname{Hom}(\mathcal{L}|U,\mathcal{O}_X|U)
  \quad\text{et}\quad
  t\in\Gamma(U,\mathcal{F}),
\]
on fait correspondre l'élément de
$\operatorname{Hom}(\mathcal{L}|U,\mathcal{F}|U)$ qui, pour tout $x\in U$,
fait correspondre à $s_x\in\mathcal{L}_x$ l'élément $u_x(s_x)t_x$ de
$\mathcal{F}_x$. Si $\mathcal{L}$ est \emph{localement libre de rang fini},
cet homomorphisme est \emph{bijectif}~; la propriété étant locale, on peut en
effet se borner au cas où $\mathcal{L}=\mathcal{O}_X^n$~; comme pour tout
$\mathcal{O}_X$-Module $\mathcal{G}$,
$\mathcal{H}\!om_{\mathcal{O}_X}(\mathcal{O}_X^n,\mathcal{G})$ est
canoniquement isomorphe à $\mathcal{G}^n$, on est ramené au cas
$\mathcal{L}=\mathcal{O}_X$, qui est immédiat.
\end{env}

\begin{env}[5.4.3]
\EGAReaderTarget{canon-ref.0.5.4.3-fr}{0.5.4.3-fr}
Si $\mathcal{L}$ est inversible, il en est de même de son dual
$\check{\mathcal{L}}=
\mathcal{H}\!om_{\mathcal{O}_X}(\mathcal{L},\mathcal{O}_X)$, car on se ramène
aussitôt (la question étant locale) au cas $\mathcal{L}=\mathcal{O}_X$. En
outre, on a un isomorphisme canonique
\[
  \mathcal{H}\!om_{\mathcal{O}_X}(\mathcal{L},\mathcal{O}_X)
  \otimes_{\mathcal{O}_X}\mathcal{L}
  \xrightarrow{\sim}
  \mathcal{O}_X
  \tag{5.4.3.1}\label{0.5.4.3.1-fr}
\]
en effet, d'après (5.4.2), il suffit de définir un isomorphisme canonique
$\mathcal{H}\!om_{\mathcal{O}_X}(\mathcal{L},\mathcal{L})
\xrightarrow{\sim}\mathcal{O}_X$. Or, pour \emph{tout}
$\mathcal{O}_X$-Module $\mathcal{F}$, on a un homomorphisme canonique
$\mathcal{O}_X\xrightarrow{\sim}
\mathcal{H}\!om_{\mathcal{O}_X}(\mathcal{F},\mathcal{F})$ (5.3.7). Il reste
à prouver que si $\mathcal{F}=\mathcal{L}$ est inversible, cet homomorphisme
est bijectif, et comme la question est locale, on est ramené au cas
$\mathcal{L}=\mathcal{O}_X$, qui est immédiat.

En raison de ce qui précède, on pose
$\mathcal{L}^{-1}=
\mathcal{H}\!om_{\mathcal{O}_X}(\mathcal{L},\mathcal{O}_X)$, et on dit que
$\mathcal{L}^{-1}$ est l'\emph{inverse} de $\mathcal{L}$. La terminologie de
\og faisceau inversible\fg{} peut se justifier de la façon suivante lorsque $X$ est
réduit à un point et $\mathcal{O}_X$ est un anneau \emph{local} $A$ d'idéal
maximal $\mathfrak m$~; si $M$ et $M'$ sont deux $A$-modules ($M$ étant de
type fini) tels que $M\otimes_A M'$ soit isomorphe à $A$, comme
$(A/\mathfrak m)\otimes_A(M\otimes_A M')$ s'identifie à
$(M/\mathfrak mM)\otimes_{A/\mathfrak m}(M'/\mathfrak mM')$, ce dernier
produit tensoriel d'espaces vectoriels sur le corps $A/\mathfrak m$ est
isomorphe à $A/\mathfrak m$, ce qui exige que $M/\mathfrak mM$ et
$M'/\mathfrak mM'$ soient de dimension 1. Pour tout élément $z\in M$
n'appartenant pas à $\mathfrak mM$, on a donc $M=Az+\mathfrak mM$, ce qui
entraîne $M=Az$ d'après le lemme de Nakayama, $M$ étant de type fini.
D'ailleurs, comme l'annulateur de $z$ annule $M\otimes_A M'$, isomorphe à
$A$, cet annulateur est $\{0\}$, et $M$ est par suite \emph{isomorphe à $A$}.
Dans le cas général, cela montre que si $\mathcal{L}$ est un
$\mathcal{O}_X$-Module de type fini, tel qu'il existe un
$\mathcal{O}_X$-Module $\mathcal{F}$ pour lequel
$\mathcal{L}\otimes_{\mathcal{O}_X}\mathcal{F}$ soit isomorphe à
$\mathcal{O}_X$, et si en outre les anneaux $\mathcal{O}_x$ sont des anneaux
locaux, alors $\mathcal{L}_x$ est un $\mathcal{O}_x$-module isomorphe à
$\mathcal{O}_x$ pour tout $x\in X$. Si $\mathcal{O}_X$ et $\mathcal{L}$ sont
supposés \emph{cohérents}, on en conclut donc que $\mathcal{L}$ est inversible
en vertu de (5.2.7).
\end{env}

\begin{env}[5.4.4]
\label{0.5.4.4-fr}
Si $\mathcal{L}$ et $\mathcal{L}'$ sont deux $\mathcal{O}_X$-Modules
inversibles, il en est de même de
$\mathcal{L}\otimes_{\mathcal{O}_X}\mathcal{L}'$, car la question étant
locale, on peut supposer que $\mathcal{L}=\mathcal{O}_X$, et le résultat est
alors trivial. Pour tout entier $n\geq1$, on désigne par
$\mathcal{L}^{\otimes n}$ le produit tensoriel de $n$ faisceaux

